\documentclass[12pt, a4paper]{article}

\usepackage[utf8]{inputenc}
\usepackage[T1]{fontenc}
\usepackage{amsmath, amssymb, amsthm}
\usepackage{mathtools}
\usepackage{geometry}
\usepackage{hyperref}
\usepackage{booktabs}
\usepackage{array}
\usepackage{enumitem}
\usepackage{algorithm}
\usepackage{algpseudocode}
\usepackage{lmodern}

\geometry{margin=2.5cm}

\hypersetup{
    colorlinks=true,
    linkcolor=blue,
    urlcolor=blue,
    citecolor=blue
}

\theoremstyle{definition}
\newtheorem{definition}{Definition}[section]
\newtheorem{remark}{Remark}[section]

\title{%
  \textbf{Mathematical Formulation of the Optimization Problem}\\[0.5em]
  \large DeliveryIQ --- Route Optimizer
}
\author{}
\date{}

\begin{document}

\maketitle
\tableofcontents
\newpage

% ─────────────────────────────────────────────────────────────────────────────
\section{Overview}
% ─────────────────────────────────────────────────────────────────────────────

DeliveryIQ solves two related combinatorial optimization problems depending on
the number of vehicles in the fleet:

\begin{enumerate}
  \item \textbf{Single-vehicle mode} (fleet size $= 1$):
        the \emph{Travelling Salesman Problem} (TSP) --- find the shortest
        round-trip tour visiting all delivery stops exactly once and returning
        to the depot.

  \item \textbf{Multi-vehicle mode} (fleet size $\geq 2$):
        a \emph{Capacitated Vehicle Routing Problem} (CVRP) --- partition
        stops across a heterogeneous fleet of vehicles with per-vehicle
        transport mode and stop-capacity constraints, then solve a TSP for each
        vehicle's assigned cluster.
\end{enumerate}

Both problems operate on a \emph{weighted directed graph} extracted from
OpenStreetMap (OSM) via the OSMnx library.  Travel times along graph edges
encode modal accessibility rules (drive / bike / walk) and serve as the
optimization objective.

% ─────────────────────────────────────────────────────────────────────────────
\section{Road-Network Graph Model}
% ─────────────────────────────────────────────────────────────────────────────

\subsection{Directed Weighted Graph}

The street network is modelled as a directed multigraph

\[
  G = (V,\, E,\, w),
\]

\noindent where:
\begin{itemize}
  \item $V$ is the set of OSM \emph{nodes} (road intersections and endpoints),
  \item $E \subseteq V \times V$ is the set of directed \emph{edges} (road
        segments),
  \item $w : E \to \mathbb{R}_{>0} \cup \{\mathcal{P}\}$ is the
        \emph{travel-time weight function}, mapping each directed edge to a
        travel time in seconds or to the penalty sentinel
        $\mathcal{P} = 10^9$ (impassable edge).
\end{itemize}

\subsection{Largest Strongly Connected Component (LSCC)}

After download, the graph is pruned to its \emph{Largest Strongly Connected
Component}:

\[
  G^* = \operatorname{LSCC}(G),
\]

\noindent ensuring that for every pair of nodes $u, v \in V^*$ a directed
path from $u$ to $v$ exists within $G^*$.  All subsequent computations use
$G^*$ exclusively.

\subsection{Modal Travel-Time Stamping}
\label{sec:modal}

Three independent copies of $G^*$ are built, one per transport mode
$m \in \mathcal{M} = \{\texttt{drive},\, \texttt{bike},\, \texttt{walk}\}$.
The travel-time weight of edge $(u, v)$ under mode $m$ is:

\begin{equation}
  \label{eq:travel_time}
  w_m(u,v) =
  \begin{cases}
    \dfrac{\ell(u,v)}{s_m}
      & \text{if mode $m$ may legally traverse $(u,v)$,} \\[8pt]
    \dfrac{\ell(u,v)}{s_m \cdot \alpha_{\text{steps}}}
      & \text{if } m = \texttt{walk} \text{ and } \texttt{highway} = \texttt{steps},\\[8pt]
    \mathcal{P}
      & \text{otherwise (mode-impassable edge),}
  \end{cases}
\end{equation}

\noindent where:
\begin{itemize}
  \item $\ell(u,v) \geq 1$ is the physical length of the edge in metres,
  \item $s_m$ is the assumed constant travel speed in metres per second
        (Table~\ref{tab:speeds}),
  \item $\alpha_{\text{steps}} = 0.5$ is the speed-reduction factor for stairs.
\end{itemize}

\begin{table}[h]
  \centering
  \caption{Assumed travel speeds by mode.}
  \label{tab:speeds}
  \begin{tabular}{lcc}
    \toprule
    Mode & Speed (km/h) & Speed (m/s) \\
    \midrule
    Drive (\texttt{drive}) & 30.0 & $8.\overline{3}$ \\
    Bicycle (\texttt{bike}) & 15.0 & $4.1\overline{6}$ \\
    Walking (\texttt{walk}) & 5.0  & $1.3\overline{8}$ \\
    \bottomrule
  \end{tabular}
\end{table}

The modal accessibility rules applied in~\eqref{eq:travel_time} are
summarised in Table~\ref{tab:modal_rules}.

\begin{table}[h]
  \centering
  \caption{Edge accessibility rules (PENALTY = impassable, else passable).}
  \label{tab:modal_rules}
  \small
  \begin{tabular}{lp{4.8cm}p{4.8cm}p{3.5cm}}
    \toprule
    Edge type & \texttt{drive} & \texttt{bike} & \texttt{walk} \\
    \midrule
    \texttt{highway} $\in$ \{footway, pedestrian, steps, \ldots\}
      & $\mathcal{P}$ & passable & passable (steps: $\times 0.5$) \\
    \texttt{motor\_vehicle} $\in$ \{no, private, destination\}
      & $\mathcal{P}$ & passable & passable \\
    \texttt{bicycle} $\in$ \{no, dismount\}
      & n/a & $\mathcal{P}$ & passable \\
    \texttt{highway} = cycleway
      & passable & passable & passable \\
    One-way reversed edge (no bike contra-flow tag)
      & --- & $\mathcal{P}$ & passable \\
    \texttt{foot} $\in$ \{no, private\}
      & n/a & n/a & $\mathcal{P}$ \\
    \bottomrule
  \end{tabular}
\end{table}

% ─────────────────────────────────────────────────────────────────────────────
\section{Shortest-Path Distance Matrix}
% ─────────────────────────────────────────────────────────────────────────────

\subsection{All-Pairs Shortest Paths (Dijkstra)}

Given the modal graph $G_m = (V^*, E, w_m)$ and a set of locations
$L = \{0, 1, \ldots, n\}$ where node $0$ is the \emph{depot} and nodes
$1, \ldots, n$ are delivery stops, the \emph{distance matrix} is defined as:

\begin{equation}
  \label{eq:matrix}
  d_m(i, j) =
  \begin{cases}
    0 & \text{if } i = j, \\
    \displaystyle\min_{\text{path } p: i \to j}
      \sum_{(u,v) \in p} w_m(u,v)
      & \text{if a finite-cost path exists,} \\
    \mathcal{P} & \text{otherwise.}
  \end{cases}
\end{equation}

The minimisation in~\eqref{eq:matrix} is computed by Dijkstra's algorithm for
every ordered pair $(i, j)$ with $i \neq j$:

\[
  d_m(i, j) = \operatorname{Dijkstra}(G_m,\, i,\, j,\, w_m).
\]

Building the full matrix requires $n(n+1)$ Dijkstra calls
($(n+1)$ nodes, $n(n+1)$ ordered pairs).

\subsection{Penalty Sentinel}

Any pair $(i, j)$ with $d_m(i, j) \geq \mathcal{P}$ is considered
\emph{mode-unreachable}.  The value $\mathcal{P} = 10^9$ seconds
($\approx 31.7$ years) is chosen to dominate any realistic travel time while
remaining a finite float, so TSP solvers can still construct a valid (if
suboptimal) tour when unreachable stops are present.

\subsection{Hybrid Last-Meter (Drive Mode)}

For drive-mode routing, a stop $j$ may lie in a pedestrian zone unreachable
by car.  In this case the car parks at the nearest car-accessible node
$v_j^{\text{car}} \in V^*$ and the driver walks the remaining distance.
The adjusted drive cost to stop $j$ from stop $i$ is:

\begin{equation}
  \label{eq:hybrid}
  d_{\text{hybrid}}(i, j) = d_{\text{drive}}(i,\, v_j^{\text{car}})
    + \frac{\ell(v_j^{\text{car}},\, j)}{s_{\text{walk}}},
\end{equation}

\noindent where $\ell(v_j^{\text{car}}, j)$ is the straight-line (Haversine)
distance from the parking node to the actual stop location.

% ─────────────────────────────────────────────────────────────────────────────
\section{Single-Vehicle: Travelling Salesman Problem}
% ─────────────────────────────────────────────────────────────────────────────

\subsection{Problem Formulation}

Let $L = \{0, 1, \ldots, n\}$ with depot $0$.  A \emph{tour} is a permutation
$\pi$ of $\{1, \ldots, n\}$ that specifies the visit order.  The total
round-trip cost is:

\begin{equation}
  \label{eq:tsp_cost}
  C(\pi) = d(0,\, \pi_1)
    + \sum_{k=1}^{n-1} d(\pi_k,\, \pi_{k+1})
    + d(\pi_n,\, 0),
\end{equation}

\noindent where $d = d_m$ is the distance matrix for the chosen mode $m$.

The TSP seeks a permutation $\pi^*$ minimising the tour cost:

\begin{equation}
  \label{eq:tsp}
  \pi^* = \underset{\pi \in \Pi_n}{\arg\min}\; C(\pi),
\end{equation}

\noindent where $\Pi_n$ is the set of all permutations of $n$ stops.

The tour is encoded as a \emph{closed sequence} starting and ending at the
depot:
\[
  [0,\; \pi_1,\; \pi_2,\; \ldots,\; \pi_n,\; 0].
\]

\subsection{Solution Methods}

The TSP is NP-hard in general.  The system selects a heuristic or
approximation algorithm automatically based on the number of stops $n$
(Table~\ref{tab:tsp_methods}).

\begin{table}[h]
  \centering
  \caption{Automatic TSP method selection.}
  \label{tab:tsp_methods}
  \begin{tabular}{cc}
    \toprule
    Stop count $n$ & Method \\
    \midrule
    $1$--$2$  & Nearest Neighbour (NN) \\
    $3$--$20$ & NN + 2-opt \\
    $21+$     & Genetic Algorithm (OX) \\
    any       & Christofides (explicit) \\
    \bottomrule
  \end{tabular}
\end{table}

\subsubsection{Nearest Neighbour (NN)}

Starting from the depot, greedily append the closest unvisited stop at each
step:

\begin{equation}
  \label{eq:nn}
  \pi_{k+1} = \underset{j \,\notin\, \{\pi_1, \ldots, \pi_k\}}{\arg\min}\;
    d(\pi_k,\, j), \qquad \pi_0 = 0.
\end{equation}

The NN heuristic runs in $O(n^2)$ time and produces a feasible tour even
when some matrix entries carry the penalty $\mathcal{P}$.

\subsubsection{2-opt Local Search}

Starting from the NN tour, iteratively swap two non-adjacent edges
$(r_i, r_{i+1})$ and $(r_j, r_{j+1})$ if the swap decreases tour cost.
For a tour $r = [r_0, r_1, \ldots, r_n, r_0]$, the \emph{2-opt swap}
of positions $i$ and $j$ reverses the sub-sequence $r_i, \ldots, r_j$:

\[
  r' = [r_0, \ldots, r_{i-1},\; r_j, r_{j-1}, \ldots, r_i,\; r_{j+1}, \ldots, r_n, r_0].
\]

The swap is accepted when:

\begin{equation}
  \label{eq:2opt}
  C(r') < C(r).
\end{equation}

The procedure repeats until no improving swap exists or a maximum of
$2\,000$ iterations is reached.  Complexity per iteration pass: $O(n^2)$.

\subsubsection{Christofides Approximation}

The Christofides algorithm gives a $\tfrac{3}{2}$-approximation for the metric
TSP.  It requires a \emph{complete, connected, undirected} helper graph
$H = (L, E_H)$ where $E_H$ contains only edges with $d(i,j) < \mathcal{P}$.
If $H$ is disconnected (i.e.\ some stop is unreachable), the algorithm falls
back to 2-opt.  When applicable, it proceeds as follows:

\begin{enumerate}[noitemsep]
  \item Compute a minimum spanning tree $T$ of $H$.
  \item Find the set $O \subset L$ of nodes with odd degree in $T$.
  \item Compute a minimum-weight perfect matching $M$ on the subgraph
        induced by $O$.
  \item Form an Eulerian multigraph $H' = T \cup M$ and extract an
        Eulerian circuit.
  \item Shortcut repeated vertices to obtain a Hamiltonian tour.
\end{enumerate}

\subsubsection{Genetic Algorithm (Order-Crossover, OX)}

For large instances ($n \geq 21$), an evolutionary algorithm is used.  The
\emph{chromosome} is a permutation $\pi$ of stops $\{1,\ldots,n\}$ (depot
excluded).  The \emph{fitness} function is:

\begin{equation}
  \label{eq:fitness}
  f(\pi) = \frac{1}{C(\pi) + 1},
\end{equation}

\noindent so lower tour cost yields higher fitness.

\paragraph{Order crossover (OX).}
Given parents $p_1$ and $p_2$, select a random sub-sequence slice $[a, b)$:

\begin{enumerate}[noitemsep]
  \item Copy $p_1[a:b]$ directly into the child at positions $[a, b)$.
  \item Fill remaining positions left-to-right with genes from $p_2$ in their
        original order, skipping genes already placed.
\end{enumerate}

\paragraph{Swap mutation.}
With probability $\mu = 0.02$, exchange two randomly chosen genes:
$\pi_i \leftrightarrow \pi_j$.

\paragraph{Generational loop.}
With population size $N_{\text{pop}} = 120$ and $G_{\text{max}} = 400$
generations:
\begin{enumerate}[noitemsep]
  \item Sort population by fitness (descending).
  \item Retain the top $10\%$ as elites.
  \item Fill the rest via OX crossover of randomly chosen elite parents,
        followed by swap mutation.
\end{enumerate}

The best chromosome after $G_{\text{max}}$ generations is returned as the
approximate solution.

% ─────────────────────────────────────────────────────────────────────────────
\section{Multi-Vehicle: Capacitated Vehicle Routing Problem}
% ─────────────────────────────────────────────────────────────────────────────

\subsection{Problem Formulation}

Let:
\begin{itemize}
  \item $S = \{s_1, \ldots, s_n\}$ be the set of delivery stops,
  \item $\mathcal{V} = \{v_1, \ldots, v_K\}$ be the fleet of $K$ vehicles,
  \item $m_k \in \mathcal{M}$ be the transport mode of vehicle $v_k$,
  \item $Q_k \in \mathbb{Z}_{>0}$ be the stop capacity of vehicle $v_k$,
  \item $d_0 \in V^*$ be the depot node.
\end{itemize}

The CVRP seeks an assignment $\sigma : S \to \mathcal{V}$ and, for each
vehicle $v_k$, a tour $\pi^{(k)}$ over the stops assigned to it, minimising
the total fleet travel time:

\begin{equation}
  \label{eq:cvrp}
  \min_{\sigma,\, \{\pi^{(k)}\}} \sum_{k=1}^{K} C_k\!\left(\pi^{(k)}\right)
\end{equation}

\noindent subject to:

\begin{align}
  \sigma(s_i) \neq \varnothing
    &\implies d_{m_{\sigma(s_i)}}(d_0, s_i) < \mathcal{P}
    & &\text{(mode-compatibility)}, \label{eq:cvrp_mode}\\
  \left|\sigma^{-1}(v_k)\right| &\leq Q_k
    & &\forall\, k \in \{1,\ldots,K\}\quad\text{(capacity)},\label{eq:cvrp_cap}\\
  \sigma^{-1}(v_k) &\cap \sigma^{-1}(v_j) = \varnothing
    & &\forall\, k \neq j\quad\text{(each stop assigned once)}, \label{eq:cvrp_part}
\end{align}

\noindent where $C_k(\pi^{(k)})$ is the tour cost \eqref{eq:tsp_cost} for
vehicle $v_k$ using distance matrix $d_{m_k}$.

Stops unreachable by any fleet vehicle ($d_m(d_0, s_i) \geq \mathcal{P}$
for all $m \in \{m_k\}$) are excluded from the optimization and reported as
warnings.

\subsection{Two-Phase Solution Algorithm}

The CVRP is solved via a \emph{cluster-first, route-second} decomposition in
two phases.

\subsubsection{Phase 1: Mode-First Stop Assignment}

Each stop $s_i$ is assigned to the transport mode $m^*(s_i)$ whose vehicles
currently have the most remaining total capacity among compatible modes:

\begin{equation}
  \label{eq:mode_assign}
  m^*(s_i) = \underset{m \in \mathcal{M}(s_i)}{\arg\max}\;
    R_m, \qquad
  R_m = \sum_{\substack{k:\, m_k = m}} \bigl(Q_k - |\sigma^{-1}(v_k)|\bigr),
\end{equation}

\noindent where $\mathcal{M}(s_i) = \{m \in \mathcal{M} : d_m(d_0, s_i) < \mathcal{P}\}$
is the set of modes compatible with stop $s_i$, and $R_m$ is the total
remaining capacity across all vehicles of mode $m$.

After mode assignment, the stops are partitioned into mode-specific pools:
$\mathcal{S}_m = \{s_i : m^*(s_i) = m\}$.

\subsubsection{Phase 2: Intra-Mode Distribution via k-Means}

Within each mode pool $\mathcal{S}_m$, stops are distributed across the
$K_m = |\{k : m_k = m\}|$ vehicles of that mode using geographic k-means
clustering.

\paragraph{k-Means clustering.}
Let $K_m' = \min(K_m,\, |\mathcal{S}_m|)$.  Run k-means on the geographic
coordinates of the stops in $\mathcal{S}_m$:

\begin{equation}
  \label{eq:kmeans}
  \min_{\{C_1,\ldots,C_{K_m'}\}}
    \sum_{j=1}^{K_m'} \sum_{s_i \in C_j}
    \left\| \begin{pmatrix} \varphi(s_i) \\ \lambda(s_i) \end{pmatrix}
           - \boldsymbol{\mu}_j \right\|^2,
\end{equation}

\noindent where $\varphi(s_i)$ and $\lambda(s_i)$ are the latitude and
longitude of stop $s_i$, $\boldsymbol{\mu}_j$ is the centroid of cluster $j$,
and $\|\cdot\|$ denotes the Euclidean norm.

\paragraph{Bearing-based cluster-to-vehicle assignment.}
To assign clusters to vehicles in a geographically coherent way, cluster
centroids are sorted by their \emph{compass bearing} from the depot.  The
bearing from point $(\varphi_1, \lambda_1)$ to $(\varphi_2, \lambda_2)$ is:

\begin{equation}
  \label{eq:bearing}
  \theta = \left(
    \arctan2\!\left(
      \sin(\Delta\lambda)\cos\varphi_2,\;
      \cos\varphi_1 \sin\varphi_2 - \sin\varphi_1 \cos\varphi_2 \cos(\Delta\lambda)
    \right) \cdot \frac{180}{\pi} + 360
  \right) \bmod 360,
\end{equation}

\noindent where $\Delta\lambda = \lambda_2 - \lambda_1$.  Clusters are
assigned to vehicles in ascending bearing order.

\paragraph{Capacity rebalancing.}
k-means ignores vehicle capacities.  After the initial assignment, any
vehicle $v_k$ that exceeds its capacity $Q_k$ transfers excess stops to the
next vehicle with remaining capacity:

\begin{equation}
  \label{eq:rebalance}
  \text{while } |\sigma^{-1}(v_k)| > Q_k :
  \quad\text{pop last stop from } \sigma^{-1}(v_k),
  \quad\text{insert into first } v_j \text{ with } |\sigma^{-1}(v_j)| < Q_j.
\end{equation}

Candidate vehicles are scanned first forward ($j > k$) then backward ($j < k$)
to avoid false over-capacity errors on the last vehicle in bearing order.
If no vehicle has remaining capacity a \texttt{ValueError} is raised.

\subsubsection{Per-Vehicle TSP}

Once the assignment $\sigma$ is determined, each vehicle $v_k$ with
$|\sigma^{-1}(v_k)| \geq 1$ independently solves the TSP
\eqref{eq:tsp}--\eqref{eq:tsp_cost} on its cluster using the modal distance
matrix $d_{m_k}$ and the auto-selected method from Table~\ref{tab:tsp_methods}.

The total fleet cost decomposes as:

\begin{equation}
  \label{eq:fleet_cost}
  C_{\text{fleet}} = \sum_{k=1}^{K} C_k(\pi^{(k)}).
\end{equation}

% ─────────────────────────────────────────────────────────────────────────────
\section{Route Reconstruction}
% ─────────────────────────────────────────────────────────────────────────────

The TSP solution $\pi^*$ is a sequence of high-level location nodes
$[0, \pi_1^*, \ldots, \pi_n^*, 0]$.  To render the route on the map, each
consecutive pair $(r_k, r_{k+1})$ is expanded to the full sequence of
intermediate OSM road nodes via Dijkstra:

\begin{equation}
  \label{eq:reconstruct}
  \text{path}(r_k, r_{k+1}) = \underset{p:\, r_k \to r_{k+1}}{\arg\min}
    \sum_{(u,v) \in p} w_m(u, v).
\end{equation}

The concatenation of all expanded sub-paths (with shared boundary nodes
deduplicated) forms the \emph{full route polyline} drawn on the Folium map.

% ─────────────────────────────────────────────────────────────────────────────
\section{Summary of Notation}
% ─────────────────────────────────────────────────────────────────────────────

\begin{table}[h]
  \centering
  \caption{Summary of mathematical notation.}
  \begin{tabular}{cl}
    \toprule
    Symbol & Description \\
    \midrule
    $G = (V, E, w)$ & Directed weighted multigraph of the road network \\
    $G^*$ & Largest Strongly Connected Component of $G$ \\
    $G_m$ & Modal copy of $G^*$ for mode $m$ \\
    $w_m(u, v)$ & Travel-time weight of edge $(u, v)$ under mode $m$ (seconds) \\
    $\mathcal{P} = 10^9$ & Impassable / unreachable sentinel value (seconds) \\
    $\ell(u, v)$ & Physical length of edge $(u, v)$ in metres \\
    $s_m$ & Travel speed for mode $m$ in metres per second \\
    $d_m(i, j)$ & Shortest-path travel time from $i$ to $j$ under mode $m$ \\
    $n$ & Number of delivery stops (excluding depot) \\
    $L = \{0, 1, \ldots, n\}$ & Set of locations; node $0$ is the depot \\
    $\pi$ & Permutation of stops (TSP tour, excluding depot) \\
    $C(\pi)$ & Total round-trip tour cost \\
    $\pi^*$ & Optimal (or near-optimal) tour \\
    $K$ & Number of vehicles in the fleet \\
    $\mathcal{V} = \{v_1, \ldots, v_K\}$ & Fleet of vehicles \\
    $m_k$ & Transport mode of vehicle $v_k$ \\
    $Q_k$ & Stop capacity of vehicle $v_k$ \\
    $\sigma : S \to \mathcal{V}$ & Stop-to-vehicle assignment function \\
    $\mathcal{S}_m$ & Set of stops assigned to mode pool $m$ \\
    $C_j$ & $j$-th k-means cluster of stops \\
    $\boldsymbol{\mu}_j$ & Centroid of cluster $C_j$ \\
    $\theta$ & Compass bearing (degrees, $[0, 360)$) \\
    $f(\pi)$ & Genetic algorithm fitness (reciprocal cost) \\
    \bottomrule
  \end{tabular}
\end{table}

\end{document}
